\section{Lorentzian Completion of the Spectral Entropy Functional}
  \label{sec:lorentzian_completion}

  \subsection{From Elliptic to Hyperbolic Dynamics}
    \label{subsec:from-elliptic-to-hyperbolic-dynamics}

    In the Riemannian formulation developed previously, the spectral entropy
    functional is defined by
    \[
      S_\Pi[g] = \frac12 \log \det{}' A_g,
    \]
    where $A_g$ is a Laplace-type elliptic operator.
    Its renormalized first variation generates an infrared Einstein response,
    with higher-derivative corrections governed by the Seeley--DeWitt
    coefficients.

    While this formulation captures the static and quasi-static structure
    of the emergent geometry, it does not address causal propagation.
    In particular, $A_g$ is elliptic and therefore encodes spatial
    correlation structure rather than hyperbolic time evolution.
    To describe gravitational waves and dynamical perturbations,
    a Lorentzian completion is required.

    We therefore introduce a hyperbolic operator $\mathcal{D}_g$ such that its spatial restriction on constant-time
    slices reduces to $A_g$.
    In a $3+1$ decomposition with lapse $N$ and shift $N^i$, the d'Alembertian reads
    \[
      \Box_g = -\frac{1}{N^2} (\partial_t - \mathcal{L}_{\vec N})^2 + \Delta_\gamma + \text{(connection terms)},
    \]
    where $\Delta_\gamma$ is the Laplacian on spatial slices.
    In the quasi-static regime $N \simeq 1$, $N^i \simeq 0$, and for slowly varying extrinsic curvature, one recovers
    \[
      \Box_g \;\longrightarrow\; -\partial_t^2 + \Delta_\gamma.
    \]
    The elliptic operator $A_g$ of the Riemannian theory is thus identified
    with the spatial block of the Lorentzian operator.

  \subsection{Definition of the Lorentzian Spectral Action}
    \label{subsec:definition-of-the-lorentzian-spectral-action}

    The Lorentzian completion of the spectral entropy functional is defined as
    \[
      S_\Pi^{(L)}[g] = \frac12 \Re \log \det{}' \mathcal{D}_g.
    \]
    Equivalently, one may formulate the theory using a Schwinger--Keldysh contour prescription,
    so that the inverse operator entering functional variations is the retarded Green operator $G_R$.

    This choice ensures:

    \begin{itemize}
      \item Reality of the quadratic kernel.
      \item Causal propagation, i.e.\ $G_R(x,y)=0$ for $x \notin J^+(y)$.
      \item Consistency with classical gravitational-wave dynamics.
    \end{itemize}

    The operator $\mathcal{D}_g$ is assumed to be globally hyperbolic and
    normally hyperbolic in the sense of Lorentzian geometry, so that
    retarded and advanced Green operators exist and are unique.

  \subsection{Second Variation and Quadratic Operator}
    \label{subsec:second-variation-and-quadratic-operator}

    The second variation of the Lorentzian spectral action takes the operatorial form
    \[
      \delta^2 S_\Pi^{(L)} = \frac12 \mathrm{Tr} \left( G_R \,\delta^2 \mathcal{D}_g \right)
      -
      \frac12 \mathrm{Tr} \left(
        G_R \,\delta \mathcal{D}_g\,
        G_R \,\delta \mathcal{D}_g
      \right).
    \]

    This expression is the hyperbolic analogue of the elliptic formula derived previously.
    The replacement $A_g^{-1} \to G_R$ is the essential modification ensuring causal structure.

    The existence and uniqueness of retarded Green operators for normally
    hyperbolic operators on globally hyperbolic spacetimes are standard
    results of Lorentzian analysis~\cite{Baer2015}.
    The use of a Schwinger--Keldysh prescription ensures a causal and
    real effective action in curved spacetime~\cite{Jordan1986,Calzetta2008}.

    Expanding about Minkowski spacetime,
    \[
      g_{\mu\nu} = \eta_{\mu\nu} + h_{\mu\nu},
    \]
    one obtains a quadratic operator of the schematic form
    \[
      \mathcal{O} = c_{\mathrm{EH}} \Box + \beta \Box^2 + \text{(non-local terms)} + \mathcal{O}(R\ell_\chi^2).
    \]

    In the infrared regime $R\ell_\chi^2 \ll 1$, the dominant local contributions are the Einstein term and the
    Weyl-squared (Bach) correction.
    Non-local contributions are suppressed by additional powers of
    $\ell_\chi^2$ and are analytic in $\Box$ at low frequency.

  \subsection{Infrared Consistency}
    \label{subsec:infrared-consistency}

    The structure above guarantees that:

    \begin{enumerate}
      \item The leading kinetic operator is hyperbolic and luminal.
      \item Higher-derivative corrections introduce an additional pole at $k^2 \sim \ell_\chi^{-2}$.
      \item In the domain $\omega \ll \ell_\chi^{-1}$,
      the additional excitation decouples,
      leaving a standard massless graviton sector.
    \end{enumerate}

    The Lorentzian completion therefore preserves the infrared Einstein
    limit while embedding it in a fully causal framework.
    Subsequent sections analyze the propagating degrees of freedom and
    derive the resulting modified dispersion relation.
